% Abhakseite „Das kann ich“ – erzeugt von rahmen.py aus den Aufgabenquelltexten
\clearpage
\hypertarget{abhaken}{}\einheitenkopf{Das kann ich}
{\small Hake ab, was du kannst.\par}\medskip
\abhakkopf{Das kennst du schon}
\abhak{Z1}{Ich kann Zahlen quadrieren, auch negative Zahlen und Dezimalzahlen.}
\abhak{Z2}{Ich kann Klammern zuerst rechnen und Punkt vor Strich beachten, auch mit negativen Zahlen.}
\abhak{Z3}{Ich kann eine lineare Gleichung durch Umformen lösen.}
\abhak{Z4}{Ich kann durch Einsetzen prüfen, ob eine Zahl eine Gleichung löst.}
\abhak{Z5}{Ich kann Quadratwurzeln ziehen und Näherungswerte runden.}
\abhak{Z6}{Ich kann an der Scheitelpunktform den Scheitel, die Öffnung und die Zahl der Nullstellen ablesen.}
\abhak{Z7}{Ich kann Terme ordnen und zusammenfassen.}
\abhak{Z8}{Ich kann Klammern ausmultiplizieren, binomische Formeln anwenden und $x$ ausklammern.}
\abhak{Z9}{Ich finde den Fehler beim Quadrieren einer Klammer.}
\abhakkopf{Einheit 1 · Wurzelziehen und Lösbarkeit}
\abhak{1}{Ich erkenne, ob eine Gleichung quadratisch oder linear ist.}
\abhak{2}{Ich erkenne, ob die rechte Seite positiv, null oder negativ ist.}
\abhak{3}{Ich kann eine Gleichung der Form $x^2 = \text{Zahl}$ durch Wurzelziehen lösen.}
\abhak{4}{Ich kann $x^2$ erst freistellen und dann die Wurzel ziehen.}
\abhak{5}{Ich kann eine Gleichung mit quadrierter Klammer rückwärts lösen.}
\abhak{6}{Ich kann durch Einsetzen prüfen, ob eine Zahl eine Lösung ist.}
\abhak{7}{Ich kann am Graphen ablesen, wie viele Lösungen eine Gleichung hat.}
\abhak{8}{Ich kann Gleichungen mit keiner, einer oder zwei Lösungen angeben und Aussagen dazu beurteilen.}
\abhak{9}{Ich finde den Fehler beim Wurzelziehen.}
\abhak{10}{Ich kann begründen, wie viele Lösungen eine Gleichung mit $x^2$ hat.}
\abhakkopf{Einheit 2 · Normalform und p-q-Formel}
\abhak{11}{Ich erkenne an der Form einer Gleichung, welcher Lösungsweg passt.}
\abhak{12}{Ich kann in der Normalform $x^2 + px + q = 0$ die Zahlen $p$ und $q$ mit Vorzeichen ablesen.}
\abhak{13}{Ich erkenne, ob die Zahl vor $x^2$ eins ist oder ob ich vorher teilen muss.}
\abhak{14}{Ich kann eine Gleichung in Normalform mit der p-q-Formel lösen.}
\abhak{15}{Ich kann eine Gleichung erst in Normalform bringen und dann mit der p-q-Formel lösen.}
\abhak{16}{Ich kann eine Gleichung in Normalform mit der p-q-Formel lösen – weiter.}
\abhak{17}{Ich kann Klammern zuerst auflösen und den passenden Lösungsweg wählen.}
\abhak{18}{Ich kann die Schnittpunkte einer Parabel mit einer Geraden berechnen.}
\abhak{19}{Ich finde den Fehler beim Lösen mit der p-q-Formel.}
\abhak{20}{Ich kann begründen, wann die p-q-Formel passt und wann es keine Lösung gibt.}
\abhak{21}{Ich kann eine Sachfrage mit der p-q-Formel beantworten.}
\abhakkopf{Einheit 3 · Sachaufgaben}
\abhak{22}{Ich erkenne, welche Lösung zur Frage passt.}
\abhak{23}{Ich kann ein Zahlenrätsel mit einer quadratischen Gleichung lösen.}
\abhak{24}{Ich kann eine Rechteckaufgabe mit einer quadratischen Gleichung lösen.}
\abhak{25}{Ich kann eine Rechteckaufgabe mit einer quadratischen Gleichung lösen – weiter.}
\abhak{26}{Ich finde den Fehler in einer Rechteckaufgabe.}
\abhak{27}{Ich kann begründen, wann beide Lösungen passen und wann nur eine.}
\abhakkopf{Ausblick Einheit 4 · Satz vom Nullprodukt}
\abhak{28}{Ich erkenne, ob bei einem Produkt rechts null steht.}
\abhak{29}{Ich kann eine Gleichung lösen, bei der ein Produkt null ist.}
\abhak{30}{Ich kann $x$ ausklammern und dann jeden Faktor null setzen.}
\abhak{31}{Ich finde den Fehler beim Lösen einer Produktgleichung.}
\abhak{32}{Ich kann begründen, warum ein Produkt null ist und warum man nicht durch $x$ teilen darf.}
\abhak{33}{Ich kann mit dem Nullprodukt eine Frage zu einem Wurf beantworten.}
